\chapter{Les variables}

\begin{Definition}[Variable]
    Une variable est un nom associé à un emplacement de la mémoire.
    On appelle ce nom un \motcle{identifiant}.
    C’est comme une boîte que l’on identifie par une étiquette.
    Elle a un \motcle{type} qui nous indique comment lire sa~\motcle{valeur} (le contenu de la boite).
    Elle a aussi une~\motcle{adresse} l'endroit où est stockée la boite dans la mémoire de l'ordinateur.
    Et une~\motcle{taille} en octet qui correspond au nombre de cases mémoires nécessaires pour stocker la valeur de la variable. 
\end{Definition}

\section*{L'identifiant}
On vient de le voir, l'identifiant est un nom que l'on attribue à une zone mémoire, en python l'identifiant peut être défini en lettres minuscules, en majuscules, avec les entiers de 0 à 9 et le caractère \_ (underscore tiret du 8).
Il peut commencer par une lettre, par un underscore, mais pas par un chiffre.
Il prend en compte la casse : mavariable $\ne$ MAVARIABLE.\\

\begin{Astuce}
    Il faut toujours bien choisir le nom de ses variables pour avoir le code le plus lisible possible.
    Idéalement, l'identifiant ne doit pas être trop long et doit décrire le rôle de la variable.

    Lorsque je dois utiliser un identifiant qui comporte plusieurs mots, je peux remplacer l'espace par le \_ (tiret du 8).\\
    Par convention, si ma variable est une constante, c'est-à-dire qu'elle ne changera jamais de valeur, son identifiant est mis en majuscule.
    Quelques exemples :
    \begin{itemize}
        \item Si ma variable sert de compteur, je peux l'appeler \textit{cpt}.
        \item Si elle sert de résultat, je peux l'appeler \textit{res}.
        \item Si ma variable contient le plus grand élément, je peux l'appeler \textit{plus\_grand}.
        \item Si ma variable contient la constante $\pi\approx3.14$, je peux l'appeler \textit{PI}.
    \end{itemize}
\end{Astuce}

\begin{Quizz}
    Je souhaite trouver un nom de variable pour conserver la couleur du texte sur mon site web, quel nom vous semble le plus approprié ?
    \begin{itemize}
    \renewcommand{\labelitemi}{$\square$}
    \item Couleur\_du\_texte
    \item COULEUR\_TXT
    \item COULEUR
    \end{itemize}
\end{Quizz}

\section*{La valeur}
Une variable sert à enregistrer une information qu'on appelle valeur. Cette valeur peut être retrouvée en appelant la variable par son identifiant.\\
Pour créer une variable, on lui donne un identifiant et on lui affecte une valeur grâce à l'opérateur $=$.

Par exemple :
\begin{lstlisting}[style=pythonstyle, language=Python]
a = 3.1415 # On attribut à la variable a la valeur 3.1415
\end{lstlisting}
\vspace{-0.8cm}
\begin{Comprendre}
Imaginons qu'on a enregistré une playlist de musiques dans l'espace mémoire de l'ordinateur. Chaque case du tableau est un espace mémoire qui contient une chanson. Chaque chanson est une \textbf{variable}.\\[0.5cm]
\begin{minipage}{0.49\textwidth}
\centering
    Comment accéder à ces variables ? \\
    \begin{tabular}{c c |c|}
    \cline{3-3}
        ? & $\mathbf{\rightarrow}$ & \large \\
         & &La vie qu'on mène - Ninho\\
         \cline{3-3}
        ? & $\mathbf{\rightarrow}$ & \large \\
          & & Not Afraid - Eminem\\
          \cline{3-3}
        ? & $\mathbf{\rightarrow}$ & \large \\
          & & Mwengé - M'toro Chamou\\
          \cline{3-3}
        ? & $\mathbf{\rightarrow}$ & \large \\
           & & La faille - Jul\\
           \cline{3-3}
    \end{tabular} 
\end{minipage}\hfill\begin{minipage}{0.49\textwidth}
 \centering   
 Grâce à leur identifiant !\\
    \begin{tabular}{c c |c|}
    \cline{3-3}
        titre1 & $\mathbf{\rightarrow}$ & \large \\
         & &La vie qu'on mène - Ninho\\
         \cline{3-3}
        titre2 & $\mathbf{\rightarrow}$ & \large \\
          & & Not Afraid - Eminem\\
          \cline{3-3}
        titre3 & $\mathbf{\rightarrow}$ & \large \\
          & & Mwengé - M'toro Chamou\\
          \cline{3-3}
        titre4 & $\mathbf{\rightarrow}$ & \large \\
           & & La faille - Jul\\
           \cline{3-3}
    \end{tabular} 
\end{minipage}
Ainsi, la variable \textit{titre1} a pour valeur \textit{La vie qu'on mène-Ninho}. Cette information sera accessible tout au long du programme.\\
En python, on crée la playlist comme suit :
\begin{lstlisting}[style=pythonstyle, language=Python]
titre1 = "La vie qu'on mène-Ninho"  
titre2 = "Not Afraid - Eminem  "
titre3 = "Mwengé - M'toro Chamou"  
titre4 = "La faille - Jul"  
\end{lstlisting}
On pourra lancer l'écoute du titre 1 quand on veut et même le changer plus tard :
\begin{lstlisting}
titre1 = "Rhythm & Blues - Ayra Starr"
\end{lstlisting}
Ainsi l'espace mémoire accessible par l'identifiant "titre1" contient \textit{Rhythm \& Blues - Ayra Starr}
\end{Comprendre}
\vspace{-1.4cm}
\begin{Quizz}
\begin{lstlisting}[style=pythonstyle, language=Python]
carte1 = "As de coeur"  
carte2 = "Roi de pique" 
carte3 = "Dame de carreaux"  
carte2 = "Valet de trèfle"  
\end{lstlisting}
A l'issue de ce programme, quelle est la représentation correcte de l'espace mémoire ?\\
\begin{minipage}{0.49\textwidth}
    \begin{itemize}
    \renewcommand{\labelitemi}{$\square$}
    \item     \begin{tabular}{c c |c|}
    \cline{3-3}
        carte1 & $\mathbf{\rightarrow}$ &  As de coeur\\
        \cline{3-3}
        carte2 & $\mathbf{\rightarrow}$ &  Roi de pique\\
        \cline{3-3}
        carte3 & $\mathbf{\rightarrow}$ & Dame de carreaux\\
        \cline{3-3}
    \end{tabular} \\[0.5cm]
    \item      \begin{tabular}{c c |c|}
    \cline{3-3}
        carte1 & $\mathbf{\rightarrow}$ &  As de coeur\\
        \cline{3-3}
        carte2 & $\mathbf{\rightarrow}$ &   Roi de pique\\
        \cline{3-3}
        carte3 & $\mathbf{\rightarrow}$ &   Dame de carreaux\\
        \cline{3-3}
          carte4 & $\mathbf{\rightarrow}$ &   Valet de trèfle\\
          \cline{3-3}
    \end{tabular} 
    \end{itemize}
\end{minipage}\begin{minipage}{0.49\textwidth}
    \begin{itemize}
    \renewcommand{\labelitemi}{$\square$}
    \item     \begin{tabular}{c c |c|}
    \cline{3-3}
        carte1 & $\mathbf{\rightarrow}$ &  As de coeur\\
        \cline{3-3}
        carte4 & $\mathbf{\rightarrow}$ &   Valet de trèfle\\
        \cline{3-3}
        carte3 & $\mathbf{\rightarrow}$ &   Dame de carreaux\\
        \cline{3-3}
        carte2 & $\mathbf{\rightarrow}$ &   Roi de pique\\
        \cline{3-3}
    \end{tabular} \\[0.5cm]
    \item      \begin{tabular}{c c |c|}
    \cline{3-3}
        carte3 & $\mathbf{\rightarrow}$ &   Dame de carreaux\\
        \cline{3-3}
        carte1 & $\mathbf{\rightarrow}$ &  As de coeur\\
        \cline{3-3}
        carte2 & $\mathbf{\rightarrow}$ &   Valet de trèfle\\
        \cline{3-3}
    \end{tabular} 
    \end{itemize}
\end{minipage}

\end{Quizz}

\section*{Le type}
Une variable a également un type, le type est important, car il dit à l'ordinateur comment lire la case mémoire qui correspond à la variable.

En effet, comme on l'a vu dans le chapitre précédent, sur un ordinateur, tout est stocké avec des bits (des 0 et des 1), mais on ne lit pas un nombre positif et un texte, de la même manière le type nous permet de savoir comment lire notre suite de bits.

En python, le type est déduit de la valeur de la variable (on parle de typage dynamique).

\begin{Comprendre}
    \begin{center}{\Huge$\ne$}\end{center}
    Supposons que votre chat soit une variable, son identifiant sera son nom, par exemple "minou" mais son type sera chat.
    Comme vous savez que minou est un chat, vous pourrez vous comporter d'une certaine manière avec lui (bien différemment que si c'était un tigre).
    En informatique, c'est exactement pareil, cpt est une variable de type int et on ne la traitera pas comme une chaine de caractère.
\end{Comprendre}

Il existe plusieurs sortes de types, certains préexistent dans le langage, on parle de~\motcle{type natifs}.
Certains correspondent à des objets élémentaires, un nombre entier, un nombre réel, un texte..., on parle de~\motcle{type primitifs}.

D'autres a des structures plus complexes, on parle de~\motcle{type construits}.\\

{\color{red}\textsc{SPOILER}}: Il est même possible de créer ses propres types, nous le verrons quand nous parlerons d'objets.\\

En python, les types primitifs les plus utiles sont :
\begin{description}
    \item[int]: pour les nombres entiers relatifs. Exemple de valeur : -410, 10, 0, 51
    \item[float]: pour les nombres réels (ou presque). Exemple de valeur : -410.25, -10.32, 0.56, 5.1, 4.0
    \item[str]: pour les chaînes de caractères. Exemple de valeur : "a", "banane", "ce prof est vraiment trop bien !"
    \item[bool]: pour les \motcle{valeurs de vérités}. Deux valeurs possibles True et False respectivement vrai et faux en anglais.
\end{description}
\begin{Quizz}
    \begin{lstlisting}[style=pythonstyle, language=Python]
age = 15 
print("J'aurai", age+1, " ans l'année prochaine")\end{lstlisting}
    Quel est le type de la variable age ?
    \begin{itemize}
    \renewcommand{\labelitemi}{$\square$}
    \item int
    \item float
    \item str
    \item bool
    \end{itemize}
\end{Quizz}
\newpage
\begin{Astuce}
    En python, pour connaitre le type d'une variable, on peut utiliser la fonction type.\\
    Voyons un exemple :
    \begin{lstlisting}[style=pythonstyle, language=Python]
a = 53
print("Le type de a est :", type(a))
# Le type de a est : <class 'int'>
    \end{lstlisting}
    De la même manière, si je veux connaître l'adresse de la variable, je peux utiliser la fonction id.\\
    Voyons un exemple :
    \begin{lstlisting}[style=pythonstyle, language=Python]
a = 53
print("L'adresse memoire où est stockée la variable a est :", id(a))
# L'adresse memoire où est stocké la variable a est 139925712759056
    \end{lstlisting}
    Enfin, si je veux connaitre la taille de la variable, je peux utiliser la fonction getsizeof.\\
    Voyons un exemple :
    \begin{lstlisting}[style=pythonstyle, language=Python]
import sys
a = 53
print("La taille mémoire de la variable a est : ", sys.getsizeof(a), "octets")
# La taille mémoire de la variable a est : 28 octets
    \end{lstlisting}
    Avec tout ça, je sais que l'objet est rangé à l'adresse 139925712759056 et qu'il occupe 28 octets. Ces 28 octets ne contiennent pas que la valeur : Python y garde aussi des informations sur l'objet (son type, le nombre de références vers lui, etc.). La valeur 53 y est écrite en binaire (0011 0101), mais pas seule.

\end{Astuce}

\begin{Quizz}
    Que va afficher le script suivant :
    \begin{lstlisting}[style=pythonstyle, language=Python]
animal = "Tortue"
print(animal, sys.getsizeof(animal), type(animal), id(animal))\end{lstlisting}
    \begin{itemize}
    \renewcommand{\labelitemi}{$\square$}
    \item Tortue 55 'str' 140581148866032
    \item 'str' Tortue 55 140581148866032
    \item animal 55 'str' 140581148866032
    \item animal 140581148866032 'str' 55 
    \end{itemize}
\end{Quizz}


\vfill

\section*{Les opérateurs}

Maintenant que nous avons vu en détail la définition précise de variable, nous allons vouloir manipuler ces variables pour cela, nous aurons besoin des opérateurs.\\

\begin{Definition}[Opérateur]
    En programmation informatique, un opérateur est un outil permet d'effectuer des opérations sur des variables. 
    Le comportement d'un opérateur dépend du type des variables avec lequel il est utilisé, mais aussi de la manière dont il est utilisé.
\end{Definition}

Voyons quelques exemples d'opérateur et leur utilisation dans des cas précis : 
\begin{description}
    \item[Les opérateurs d'affectation] :
    \begin{description}
        \item[=] : Permet  d'affecter une valeur à une variable.
    \end{description}
    \item[Les opérateurs arithmétiques] :
    \begin{description}
        \item[+] : Permet d'additionner deux nombres.
        \item[-] : Permet de soustraire deux nombres.
        \item[*] : Permet de multiplier deux nombres.
        \item[/] : Permet de diviser deux nombres.
        \item[**] : Permet d'effectuer l'opération puissance.
        \item[//] : Permet de calculer le quotient de la division entiere.
        \item[\%] : Permet de calculer le reste de la division entiere.
    \end{description}
\item[Les opérateurs de comparaison] :
    \begin{description}
        \item[==] : Compare si deux éléments sont égaux.
        \item[!=] : Compare si deux éléments sont différents.
        \item[>] : Compare si l'élément de gauche est supérieur (strictement plus grand) à l'élément de droite.
        \item[<] : Compare si l'élément de gauche est inférieur (strictement plus petit) à l'élément de droite.
        \item[>=] : Compare si l'élément de gauche est supérieur (plus grand) ou égal à l'élément de droite
        \item[<=] : Compare si l'élément de gauche est inférieur (plus petit) ou égal à l'élément de droite
    \end{description}
    \item[Encore plein d'autres] que vous découvrirez au fil du temps (is, in, as, from, :=, \&, or, and, [], not, +=, *=,  -=, del, def, lambda...)
\end{description}

Sans surprise , nous pouvons utiliser les opérateurs : *, +, -, / sur les variables de type int ou float.\\

Créons un code python qui additionne des pommes et des mangues et qui affiche la somme des deux :


\begin{lstlisting}[style=pythonstyle, language=Python]
pommes = 47328 # on affecte la valeur 47328 à la variable pommes
mangues = 879443 # on affecte la valeur 879443 à la variable mangues
fruits_total = pommes + mangues # on additionne deux variables
print(fruits_total)
\end{lstlisting}
Ce programme affiche \textit{926771}. Remarquez qu'on a \textbf{pas besoin de RECOPIER} les valeurs, on peut directement utiliser \textbf{l'identifiant de la variable}.
\newpage
\begin{Quizz}
    \begin{lstlisting}[style=pythonstyle, language=Python]
nb_coco_totale = 20
coco_vertes = 5
coco_oranges = 10
coco_marrons = 3
coco_pourries = 2 \end{lstlisting}

    Quelle ligne peut-on ajouter pour afficher le nombre de cocos non pourries ?
    \begin{itemize}
    \renewcommand{\labelitemi}{$\square$}
    \item  \begin{lstlisting}[style=pythonstyle, language=Python]
print(nb_coco_totale - coco_pourries)\end{lstlisting}

    \item \begin{lstlisting}[style=pythonstyle, language=Python]
print(coco_vertes + coco_oranges + coco_marrons)\end{lstlisting}

    \item \begin{lstlisting}[style=pythonstyle, language=Python]
print((coco_vertes + coco_oranges + coco_marrons)/2)\end{lstlisting}

    \item  \begin{lstlisting}[style=pythonstyle, language=Python]
print(nb_coco_totale - (coco_vertes + coco_oranges + coco_marrons))\end{lstlisting}
    \end{itemize}
\end{Quizz}

\begin{Astuce}
    En python, on peut additionner(+) les chaînes de caractère, mais on ne peut pas les soustraire, ni les diviser, ni les multiplier. Enfin, vous pouvez essayer, mais vous aurez une erreur  !
    \begin{lstlisting}[style=pythonstyle, language=Python]
        prenom = "Alice "
        phrase = "adore l'informatique"
        print(prenom+phrase)
        # Alice adore l'informatique\end{lstlisting}
        
    On dit qu'on \textbf{CONCATÈNE} deux chaînes de caractères\\
    \begin{lstlisting}[style=pythonstyle, language=Python]
        prenom = "Alice "
        phrase = "adore l'informatique"
        print(prenom*phrase)
        # TypeError: can't multiply sequence by non-int of type 'str'
    \end{lstlisting}
    Type d'erreur : on ne peut pas multiplier une séquence de non-int avec un type 'str'.\\
    Vous découvrirez petit à petit comment manipuler ces opérateurs.
\end{Astuce}
